\section{Cyclic support of the bands}\label{sec:path-band}

For the band step of the induction, it remains to determine the support
needed in $U_{m,1}$ and $U_{m,2}$ for \Cref{prop:band-quotient}.
The width-one factor is trivial.
For width two we will detect every shape with at most four rows directly
in the specified generator, using the invariant evaluation pairing.

For width one, the band consists of the $m$ positions with $i+j=m+1$.
Every row and column block is a singleton, and both alphabets consist only
of the letter $1$. Thus $u_{m,1}$ is a nonzero constant word tensor, fixed
by every position permutation, and
\begin{equation}\label{path:one-band}
 U_{m,1}\cong S^{(m)} \qquad (m\geq2).
\end{equation}
We now determine the support needed from the particular width-two
generator.

\subsection{The indexed path contraction}

Let $m\geq3$, $K=2m-1$, and $F=B_m\setminus B_{m-2}$.
Number its positions along the path by
\begin{equation}\label{path:positions}
 p_{2j-1}=(m+1-j,j)\quad(1\leq j\leq m),\qquad
 p_{2j}=(m-j,j)\quad(1\leq j<m).
\end{equation}
In this numbering, the row blocks are the singleton $1$ and the pairs
$(2,3),(4,5),\ldots,(K-1,K)$; the column blocks are
$(1,2),(3,4),\ldots,(K-2,K-1)$ and the singleton $K$.
Order each pair by increasing path index and let $u$ denote the resulting
alternating band tensor. The row orders agree with increasing $j$, while
each of the $m-1$ column pairs reverses increasing $i$. Consequently
\begin{equation}\label{path:generator-sign}
 u=(-1)^{m-1}u_{m,2},\qquad U_{m,2}=\C[S_F]u.
\end{equation}
\Cref{fig:band-path} illustrates the numbering and the consecutive
segments used in \Cref{path:conjugation-example}.

\begin{figure}[ht]
\centering
\setlength{\unitlength}{1pt}
\begin{picture}(430,155)
\small
\put(90,148){\makebox(0,0){Band in $B_3$}}
\put(10,139){\vector(1,0){25}}
\put(39,139){\makebox(0,0){$j$}}
\put(10,139){\vector(0,-1){25}}
\put(10,108){\makebox(0,0){$i$}}
\put(45,110){\color[gray]{0.65}\circle*{6}}
\put(45,123){\makebox(0,0){$D=B_1$}}
\put(45,30){\circle*{4}}
\put(45,70){\circle*{4}}
\put(90,70){\circle*{4}}
\put(90,110){\circle*{4}}
\put(135,110){\circle*{4}}
\put(45,34){\vector(0,1){32}}
\put(90,74){\vector(0,1){32}}
\multiput(49,70)(8,0){4}{\line(1,0){4}}
\put(81,70){\vector(1,0){5}}
\multiput(94,110)(8,0){4}{\line(1,0){4}}
\put(126,110){\vector(1,0){5}}
\put(45,16){\makebox(0,0){$p_1$}}
\put(30,70){\makebox(0,0){$p_2$}}
\put(106,70){\makebox(0,0){$p_3$}}
\put(90,124){\makebox(0,0){$p_4$}}
\put(135,124){\makebox(0,0){$p_5$}}
\put(305,148){\makebox(0,0){Dual-test columns for $\eta=(2,2,1)$}}
\multiput(220,90)(40,0){3}{\circle*{4}}
\put(350,90){\circle*{4}}
\put(390,90){\circle*{4}}
\put(224,90){\vector(1,0){32}}
\multiput(264,90)(8,0){3}{\line(1,0){4}}
\put(288,90){\vector(1,0){8}}
\put(304,90){\vector(1,0){42}}
\multiput(354,90)(8,0){3}{\line(1,0){4}}
\put(378,90){\vector(1,0){8}}
\put(220,76){\makebox(0,0){$p_1$}}
\put(260,76){\makebox(0,0){$p_2$}}
\put(300,76){\makebox(0,0){$p_3$}}
\put(350,76){\makebox(0,0){$p_4$}}
\put(390,76){\makebox(0,0){$p_5$}}
\put(325,70){\color[gray]{0.5}\line(0,1){38}}
\put(215,55){\line(1,0){90}}
\put(215,55){\line(0,1){6}}
\put(305,55){\line(0,1){6}}
\put(260,43){\makebox(0,0){$r_1=3$}}
\put(335,55){\line(1,0){60}}
\put(335,55){\line(0,1){6}}
\put(395,55){\line(0,1){6}}
\put(365,43){\makebox(0,0){$r_2=2$}}
\end{picture}
\caption{The width-two band for $m=3$. Solid arrows are column pairs;
dashed arrows are row pairs, all directed by increasing path index.
Column arrows reverse increasing $i$, as in
\Cref{path:generator-sign}. The row singleton is $p_1$ and the column
singleton is $p_5$, each carrying the fixed letter $1$ in its layer.
On the right, the dual-test tableau of shape $\eta=(2,2,1)$ has column
lengths $r_1=3$ and $r_2=2$, giving the two bracketed consecutive segments.
The gray vertical line separates these segments, which differ from the
original pair blocks.}
\label{fig:band-path}
\end{figure}

Put $E=\C^2$ with basis $e_1,e_2$ and $V=E\otimes E$.
Ordinary tensor reassociation identifies the two band layers with
$V^{\otimes K}$, using the position order in \Cref{path:positions}.
Identify $V^*$ with $\Mat_2(\C)$ by letting a matrix $X$ take the value
$X_{ab}$ on $e_a\otimes e_b$. All pairings below are bilinear evaluation
pairings. Write
\[
 J=\begin{pmatrix}0&1\\-1&0\end{pmatrix}.
\]

\begin{lemma}[Path contraction]\label{path:contraction}
For arbitrary $X_1,\ldots,X_K\in\Mat_2(\C)$,
\begin{equation}\label{path:contraction-formula}
 \left\langle u,X_1\otimes\cdots\otimes X_K\right\rangle
 =(-1)^{m-1}
 \left[X_1\adj(X_2)X_3\adj(X_4)\cdots
                 X_{K-2}\adj(X_{K-1})X_K\right]_{11}.
\end{equation}
\end{lemma}

\begin{proof}
The two singleton blocks require the first row-layer index and last
column-layer index to be $1$. The coefficient of an alternating pair in
path order is $J_{ab}$. Hence the left side of
\Cref{path:contraction-formula} is exactly
\begin{equation}\label{path:index-sum}
 \sum_{\substack{a_1,\ldots,a_K,b_1,\ldots,b_K\in\{1,2\}\\
                  a_1=b_K=1}}
 \left(\prod_{\ell=1}^{m-1}
   J_{b_{2\ell-1},b_{2\ell}}J_{a_{2\ell},a_{2\ell+1}}\right)
 \prod_{q=1}^K(X_q)_{a_qb_q}.
\end{equation}
Reading the contracted indices successively from $a_1$ to $b_K$ gives
\[
 \left[X_1JX_2^{\mathsf T}JX_3JX_4^{\mathsf T}J
       \cdots X_{K-2}JX_{K-1}^{\mathsf T}JX_K\right]_{11}.
\]
The transpose at an even position records that its column index is
encountered before its row index. For
$X=\left(\begin{smallmatrix}a&b\\c&d\end{smallmatrix}\right)$,
direct multiplication gives
\begin{equation}\label{path:adjugate}
 \adj(X)=\begin{pmatrix}d&-b\\-c&a\end{pmatrix}
        =\tr(X)I-X,\qquad JX^{\mathsf T}J=-\adj(X).
\end{equation}
There are $m-1$ even positions, yielding the stated sign.
In particular, adjugation is a linear operation on $\Mat_2(\C)$.
\end{proof}

\subsection{Alternating consecutive segments}

For a shape with at most four rows, place its tableau columns on
successive disjoint segments of the path. On a segment of length $r$,
the dual polytabloid alternates the first $r$ elements of one common
basis of $V^*$. The contraction formula therefore assigns an alternating
matrix sum to each segment; these sums multiply in path order.
We now choose the basis so that every such segment matrix, for
$1\le r\le4$ and either starting parity, is invertible.

Use the following ordered basis of $V^*$:
\begin{equation}\label{path:matrix-basis}
 M_1=I,\qquad
 M_2=Z=\begin{pmatrix}1&0\\0&-1\end{pmatrix},\qquad
 M_3=T=\begin{pmatrix}0&1\\1&0\end{pmatrix},\qquad
 M_4=J=ZT.
\end{equation}
These matrices are linearly independent: $I,Z$ span the diagonal
matrices and $T,J$ span the off-diagonal matrices. The last three are
traceless and anticommute in pairs; moreover
\[
 Z^2=T^2=I,\qquad J^2=-I.
\]
Thus all four are invertible. Alternating evaluations on $Z,T,J$ are also
used to test polynomial identities for two-by-two matrices with adjugation
\cite[Section~3]{drenskygiambruno1994}; here the adjugations are determined
by the path positions.

For a global path position $q$, define the linear operation
\[
 D_q(X)=
 \begin{cases}
 X,&q\text{ odd},\\
 \adj(X),&q\text{ even}.
 \end{cases}
\]
For a segment of length $r\leq4$ beginning at position $p$, define its
alternating matrix by the following sum of ordered products:
\begin{equation}\label{path:segment-definition}
 A_{r,p}=\sum_{\pi\in S_r}\sgn(\pi)
      D_p(M_{\pi(1)})D_{p+1}(M_{\pi(2)})\cdots
      D_{p+r-1}(M_{\pi(r)}).
\end{equation}

\begin{lemma}[Segment alternation]\label{path:segment}
Let $e$ be the number of even positions among $p,\ldots,p+r-1$, and put
$\delta=1$ if $p$ is even and $\delta=0$ otherwise. Then
\begin{equation}\label{path:segment-formula}
 A_{r,p}=r!(-1)^{e-\delta}M_1\cdots M_r.
\end{equation}
In particular, all these matrices are invertible. Their values are
\begin{equation}\label{path:segment-table}
 \begin{array}{c|cc}
 r & p\text{ odd} & p\text{ even}\\ \hline
 1&I&I\\
 2&-2Z&2Z\\
 3&-6J&-6J\\
 4&-24I&24I
 \end{array}
\end{equation}
\end{lemma}

\begin{proof}
Fix a permutation $\pi$ and let $j$ be the local slot occupied by $M_1=I$.
Deleting $I$ leaves a permutation of $M_2,\ldots,M_r$.
Its sign is exactly the sign acquired when those pairwise anticommuting
matrices are reordered. Moving $I$ from its first canonical slot to slot
$j$ contributes $(-1)^{j-1}$ to $\sgn(\pi)$ and changes no matrix product.
It follows that
\[
 \sgn(\pi)M_{\pi(1)}\cdots M_{\pi(r)}
       =(-1)^{j-1}M_1\cdots M_r.
\]
By \Cref{path:adjugate}, adjugation fixes $I$ and negates each of the
other basis matrices. Write $\epsilon_j=1$ if $p+j-1$ is even, and $0$
otherwise. Adjugation therefore contributes the additional sign
$(-1)^{e-\epsilon_j}$. Since
$\epsilon_j\equiv\delta+j-1\pmod2$, the total sign is
$(-1)^{e-\delta}$, independently of $\pi$. All $r!$ summands in
\Cref{path:segment-definition} consequently agree, proving
\Cref{path:segment-formula}, including $r=1$.
Finally, the canonical products for $r=1,2,3,4$ are respectively
$I,Z,J,-I$, which gives \Cref{path:segment-table}.
\end{proof}

\subsection{Detection by a dual polytabloid}

The segment calculation gives an invertible ordered product. Its fixed
endpoint entry can nevertheless vanish. We now choose the dual test
basis so that this entry is nonzero, while keeping the band generator
and its endpoint letters unchanged.

\begin{proposition}[Four-row support of the width-two band]\label{prop:two-band}
For every $m\geq3$ and every partition $\eta\vdash 2m-1$ with
$\len(\eta)\leq4$, the irreducible module $S^\eta$ occurs in $U_{m,2}$.
\end{proposition}

\begin{proof}
Let $r_1,\ldots,r_c$ be the column lengths of $\eta$ in their usual order,
so $1\leq r_j\leq4$ and $\sum_jr_j=K$. Partition the path into consecutive
segments of these lengths, beginning at
$s_j=1+\sum_{i<j}r_i$.
For any ordered basis $N_1,\ldots,N_4$ of $V^*$, let $t_N$ be the dual
tensor whose factor on segment $j$ is
\[
 \sum_{\pi\in S_{r_j}}\sgn(\pi)
       N_{\pi(1)}\otimes\cdots\otimes N_{\pi(r_j)}.
\]
The segments are the columns of a tableau of shape $\eta$ and the initial
basis entries label its rows. By \Cref{lem:polytabloid}, their product
$t_N$ is a polytabloid, and
\[
 \mathcal T_N=\C[S_F]t_N\ \subseteq\ (V^*)^{\otimes K}
 \quad\text{is irreducible of type }\eta.
\]

First take $N_i=M_i$. Multilinearity of
\Cref{path:contraction-formula} and the definition of $t_M$ give
\begin{equation}\label{path:ordered-factorization}
 \langle u,t_M\rangle=(-1)^{m-1}[B]_{11},\qquad
 B=A_{r_1,s_1}A_{r_2,s_2}\cdots A_{r_c,s_c}.
\end{equation}
Indeed, summing independently over the permutations within each segment
is exactly the distributive expansion of the displayed ordered product.
The segments are consecutive, so no matrix factors are interchanged in
this factorization. By \Cref{path:segment}, $B$ is invertible.

The identity $\adj(X)=\tr(X)I-X$ also shows that, for every
$Q\in\GL_2(\C)$,
\[
 D_q(QXQ^{-1})=QD_q(X)Q^{-1}.
\]
Choose an eigenvector $v$ of $B$ with eigenvalue $\beta$; invertibility
gives $\beta\ne0$. Choose $Q$ taking $v$ to the first coordinate vector
and set $N_i=QM_iQ^{-1}$. Conjugation is an invertible linear operation
on $V^*$, so these $N_i$ are still an ordered basis. Each alternating
segment matrix is now $QA_{r_j,s_j}Q^{-1}$, and their ordered product is
$QBQ^{-1}$. Since $(QBQ^{-1})e_1=\beta e_1$,
\begin{equation}\label{path:nonzero-pairing}
 \langle u,t_N\rangle=(-1)^{m-1}[QBQ^{-1}]_{11}
                    =(-1)^{m-1}\beta\ne0.
\end{equation}
Only the test covectors have changed; $u$ and its endpoint letters remain
those fixed in \Cref{path:generator-sign,path:index-sum}.

The invariant evaluation pairing defines
\begin{equation}\label{path:quotient-map}
 \Phi:U_{m,2}\longrightarrow\mathcal T_N^*,\qquad
 \Phi(x)(t)=\langle x,t\rangle.
\end{equation}
For $g\in S_F$, its equivariance follows explicitly from
$\Phi(gx)(t)=\langle x,g^{-1}t\rangle=(g\Phi(x))(t)$.
By \Cref{path:nonzero-pairing} this map is nonzero on the particular
generator $u$. Its target is irreducible and, by the self-duality in
\Cref{tools:dual-sign}, has type $\eta$. Thus $\Phi$ is surjective, and
semisimplicity proves the claimed occurrence in $U_{m,2}$.
\end{proof}

\begin{example}\label{path:conjugation-example}
The change of test basis can be essential. For $m=3$ and
$\eta=(2,2,1)$, the column segments have lengths $3,2$, beginning at
positions $1,4$. Their product in \Cref{path:ordered-factorization} is
$(-6J)(2Z)=12T$, whose $(1,1)$ entry vanishes. With
\[
 Q=\begin{pmatrix}1&1\\1&2\end{pmatrix},\qquad
 QTQ^{-1}=\begin{pmatrix}1&0\\3&-1\end{pmatrix},
\]
the resulting pairing is $12$, since $(-1)^{m-1}=1$.
\end{example}
